\chapter{12.9 等腰三角形}



\section*{三、解答题（共4小题）}

\solutionlist[14]


\item 阅读下列材料，按要求解答问题：

如图 1，在 $\triangle ABC$ 中，$\angle A = 2\angle B$，且 $\angle A = 60^\circ$。小明通过以下计算：由题意，$\angle B = 30^\circ$，$\angle C = 90^\circ$，$c = 2b$，$a = \sqrt{3}b$，得 $a^2 - b^2 = (\sqrt{3}b)^2 - b^2 = 2b^2 = b \cdot c$。即 $a^2 - b^2 = bc$。于是，小明猜测：对于任意的 $\triangle ABC$，当 $\angle A = 2\angle B$ 时，关系式 $a^2 - b^2 = bc$ 都成立。

(1) 如图 2，请你用以上小明的办法，对等腰直角三角形进行验证，判断小明的猜测是否正确，并写出验证过程。

\rightquestionimage[width=0.5\textwidth]{0.5\textwidth}{12.9-14-1.png}

\vspace{3cm}

(2) 如图 3，$\angle A = 30^\circ$，$\angle B = 15^\circ$，试设计方法，判断小明的猜测是否正确，并写出验证过程。

\rightquestionimage[width=0.5\textwidth]{0.5\textwidth}{12.9-14-2.png}

\vspace{3cm}

(3) 若一个三角形的三边长恰为三个连续偶数，且 $\angle A = 2\angle B$，请直接写出这个三角形三边的长不必说明理由。

\vspace{3cm}



\item 如图，在平面直角坐标系中，已知点 $A(0, 8)$，$B(6, 0)$。以 $\triangle AOB$ 的一边为边画等腰三角形，使它的第三个顶点在 $\triangle AOB$ 的一边上。请在图①、图②中分别画出一个符合条件的等腰三角形，且两个图中的等腰三角形不相同，在图中标明所画等腰三角形的第三个顶点的坐标（不要求尺规作图，不写求解过程）。

% 插入图片：此处应有图①

% 插入图片：此处应有图②

\vspace{6cm}



\item 如图，在 $\triangle ABC$ 中，$AB = AC = a$，$BC = b$，且 $2a > b$，$BG \perp AC$ 于 $G$，$DE \perp AB$ 于 $E$，$DF \perp AC$ 于 $F$。

(1) 在图 (1) 中，$D$ 是 $BC$ 边上的中点，计算 $DE + DF$ 和 $BG$ 的长（用 $a, b$ 表示），并判断 $DE + DF$ 与 $BG$ 的关系。


(2) 在图 (2) 中，$D$ 是线段 $BC$ 上的任意一点，$DE + DF$ 与 $BG$ 的关系是否仍然成立？如果成立，证明你的结论；如果不成立，请说明理由。


(3) 在图 (3) 中，$D$ 是线段 $BC$ 延长线上的点，探究 $DE, DF$ 与 $BG$ 的关系。（不要求证明）

\centerquestionimage[width=0.7\textwidth]{0.7\textwidth}{12.9-16.png}

\vspace{5cm}


\newpage

\item 如图所示，在 $\triangle ABC$ 中，$AB = AC$，$\angle BAC = 80^\circ$，$P$ 在 $\triangle ABC$ 内，$\angle PBC = 10^\circ$，$\angle PCB = 30^\circ$，求 $\angle PAB$ 的度数。

\rightquestionimage[width=0.5\textwidth]{0.5\textwidth}{12.9-17.png}

\vspace{3cm}

\end{enumerate}
